\section{Contexto y justificación}

\begin{nota}{Marco general.}
Esta situación de aprendizaje se integra en la programación general de Matemáticas~I. Por eso este apartado describe un marco común para un instituto público de la Comunidad de Madrid, y deja para cada centro la concreción de sus datos y de su grupo. El IES Ejemplo que aparece en el reto es un centro ficticio que sirve de ilustración.
\end{nota}

\subsection{Lugar en la programación}
La propuesta se sitúa en el segundo trimestre de Matemáticas~I, cuando se abordan la geometría y la trigonometría. Reúne en una sola unidad los vectores, la trigonometría y las rectas, y se planifica para 20 sesiones de 55 minutos, una duración que cada centro puede ajustar a su calendario.

\subsection{Condiciones del centro}
La propuesta puede aplicarse en cualquier instituto con Bachillerato de la modalidad de Ciencias y Tecnología. Solo requiere edificios u objetos altos que se puedan medir desde el exterior, un espacio al aire libre seguro, ordenadores con conexión a internet y permiso para trabajar sobre los espacios del propio centro. El reto gana sentido si el centro ya dispone de paneles solares o estudia instalarlos, pero también funciona como estudio de viabilidad. Cada centro adapta los objetos que se miden, los datos de latitud y las personas que reciben las propuestas.

\subsection{El alumnado}
La propuesta se dirige a estudiantes de 1.º de Bachillerato, habitualmente de 16 y 17 años. Se parte de un grupo heterogéneo en Matemáticas, donde convive quien domina los procedimientos pero tiene dificultad para justificarlos con quien necesita apoyo para pasar del triángulo rectángulo a los problemas generales. La programación prevé medidas generales para esa diversidad y orientaciones para las necesidades más habituales (sección~\ref{sec:diversidad}), que cada docente concreta según su grupo.

\subsection{Organización y tutoría}
Matemáticas I es una materia específica de modalidad con cuatro horas lectivas semanales \parencite[anexo~I, p.~32]{decreto64}. El departamento de Matemáticas coordina la propuesta dentro de su programación, y la comisión de coordinación pedagógica la articula con el resto de materias. El profesor tutor realiza el seguimiento del proceso educativo del grupo \parencite[art.~5.1, p.~17]{decreto64}, y el profesor de Matemáticas le informa del funcionamiento de los equipos. Como el horario del anexo~I no incluye una hora de tutoría en este curso \parencite[p.~32]{decreto64}, la formación en los roles se realiza dentro de las sesiones de Matemáticas\soloProducto{ (sección~\ref{sec:cooperativo})}. El departamento de orientación puede aportar información sobre salidas profesionales vinculadas a la energía, en línea con la finalidad orientadora del plan de acción tutorial \parencite[art.~5.3, p.~17]{decreto64}.

\subsection{Documentos del centro}
La propuesta debe ser coherente con los documentos institucionales de cada centro:
\begin{itemize}
  \item \textbf{Proyecto educativo:} recoge valores, fines y prioridades del centro e incorpora la concreción del currículo \parencite[art.~121.1, p.~122921]{lomloe}. La propuesta encaja especialmente si el centro prioriza la sostenibilidad.
  \item \textbf{Programación general anual:} el Consejo Escolar la aprueba y evalúa \parencite[art.~127.b, p.~122923]{lomloe}. Puede recoger la mejora energética del centro como objetivo del curso.
  \item \textbf{Plan de acción tutorial:} forma parte de la programación general anual \parencite[art.~5.3, p.~17]{decreto64}.
  \item \textbf{Programación didáctica de Matemáticas~I:} sitúa esta propuesta dentro de la secuencia anual de la materia.
\end{itemize}

\subsection{Marco legislativo}
La propuesta se rige por la Ley Orgánica 3/2020 \parencite{lomloe}, el Real Decreto 243/2022, que fija las enseñanzas mínimas y los descriptores del perfil de salida \parencite{rd243}, y el Decreto 64/2022 de la Comunidad de Madrid \parencite{decreto64}, cuyo anexo~II recoge las competencias específicas, los criterios y los contenidos de Matemáticas~I. El decreto pide planificar actividades que pongan en contexto la enseñanza y el aprendizaje \parencite[art.~16.3, p.~21]{decreto64} y permite a los centros adoptar innovaciones pedagógicas en el marco de su autonomía \parencite[art.~19.3.a, pp.~21--22]{decreto64}.

\subsection{Justificación de la propuesta}
\textbf{Por qué esta propuesta en este curso.} El currículo destaca el carácter instrumental de las matemáticas en otras áreas científicas y tecnológicas y sitúa la resolución de problemas y la investigación como aspectos fundamentales de la materia \parencite[p.~320]{decreto64}. Vectores, trigonometría y rectas suelen enseñarse por separado y este alumnado tiende a verlos sin relación. El problema de la sombra los necesita a todos: el rayo de sol es una recta con un vector director, su ángulo procede de la trigonometría y las alturas inaccesibles se obtienen resolviendo triángulos.

\textbf{Por qué con este enfoque.} Es frecuente que el alumnado de este curso domine los procedimientos, pero le cueste justificarlos. Defender una decisión real ante quien decide obliga a argumentar, y los equipos heterogéneos aprovechan la diversidad de niveles.

\subsection{Relación con otras materias}
\begin{itemize}
  \item \textbf{Tecnología e Ingeniería I:} energías renovables, eficiencia energética y sostenibilidad \parencite[p.~385]{decreto64}. Aporta los criterios técnicos de la instalación.
  \item \textbf{Física y Química:} carácter vectorial de las magnitudes y composición vectorial \parencite[p.~166]{decreto64}.
  \item \textbf{Orientación:} salidas profesionales vinculadas a la energía.
\end{itemize}

\subsection{Objetivos de Desarrollo Sostenible}
\begin{itemize}
  \item \textbf{ODS~7} \parencite{onu7}: la meta~7.2 propone aumentar la proporción de energía renovable. El alumnado estudia cómo aprovechar mejor una instalación solar real.
  \item \textbf{ODS~11} \parencite{onu11}: la meta~11.6 busca reducir el impacto ambiental de las ciudades. El instituto actúa como pequeña comunidad que planifica el uso de su propio espacio.
\end{itemize}
